\documentclass{situationist-preprint}

\title{Priority Monism Defended}
\msid{MET-09-a}
\msfamily{MET-09}
\msdiscipline{Metaphysics}
\mstype{Principal manuscript}
\mskeywords{priority monism; fundamentality; grounding; cosmos; parts and wholes; totality facts; laws of nature; entanglement}

\begin{document}
\maketitle

\begin{abstract}
Is reality fundamentally one? We argue that it is: the cosmos is the only fundamental object, and every proper part of it, including every ordinary object, is an abstraction from it, in the sense that its existence is grounded in a fact about the cosmos. The thesis is priority monism, not existence monism. Tables and organisms exist, as MET-06 argues, and the cosmos is a composite of them and of everything else; what the thesis concerns is the direction of grounding. The argument turns on an asymmetry between the cosmos and its parts. Every proper part has an environment, and whether its parts compose it depends on the dynamics to which they are subject, which depend on the laws, which depend in turn on the whole mosaic, including the totality fact that the cosmos contains nothing more. The cosmos alone has no environment, and its parts share a common fate only vacuously, because it is all there is. Given that grounding is a strict partial order and that facts about an object are grounded in part in its existence, the cosmos depends on nothing and everything else depends on it. The result is compatible with MET-10: there is no fundamental level of parts, and the one fundamental object is the whole.
\end{abstract}

\section{The question}\label{s:question}

Monism is the thesis that reality is one. Stated so, it admits two readings, and the history of the question has been confused by the failure to separate them \citep{Schaffer2010}. \emph{Existence monism} holds that exactly one concrete object exists: there are no tables, no organisms, no electrons, only the one thing of which these are, at best, modes or appearances. \emph{Priority monism} holds that exactly one concrete object is fundamental: there are tables, organisms and electrons, but they exist and are as they are in virtue of how the one whole is, and the whole does not exist in virtue of them. Pluralism, correspondingly, divides into the denial of existence monism, which is nearly universal, and the denial of priority monism, which is the orthodoxy of analytic metaphysics since Russell: the view that the fundamental objects are many and small, and that the whole is a mere sum of them.

We defend priority monism.

\begin{quote}
\textbf{Thesis.} The cosmos is the only fundamental object. Every other concrete object, including every ordinary object, is an abstraction from it: a proper part of the cosmos whose existence is grounded in a fact about the cosmos.
\end{quote}

Three features of the thesis need emphasis at the outset. First, it is not existence monism. MET-06 argues that some pluralities compose objects, exactly those that share a common fate under the relevant dynamics, and that tables and organisms exist. We accept that result, and we accept its verdict that the cosmos is itself a composite (MET-06-a, Section~6). The thesis concerns the direction of grounding between the composite cosmos and its parts; that the cosmos has parts and that they compose it is common ground. Second, \enquote{abstraction} is a technical term, defined below, and it does not mean that ordinary objects are unreal, mental or conventional. An abstraction from the cosmos is as real as the cosmos; it is derivative. Third, the thesis is compatible with the conclusion of MET-10 that there is no fundamental level of reality. A level is a decomposition of the cosmos into proper parts, and the cosmos is not a level; MET-10 says that no level of parts is fundamental, and we say that the whole is. Section~\ref{s:levels} shows that the two theses form a single coherent picture.

The argument rests on an asymmetry between the cosmos and everything else. Every proper part of the cosmos has an environment, the rest of the cosmos. Whether some things compose an object depends, on the criterion of MET-06, on how they respond to the perturbations their environment characteristically produces under the laws, and the laws, on the best-system account of MET-12, depend on the whole mosaic, including the fact that the mosaic contains nothing more than it does. So the existence of every proper part depends on a fact about the whole. The cosmos has no environment. Its parts share a common fate, but only vacuously, because there is nothing outside them to perturb them; and that is a fact about the cosmos, not about its parts. The asymmetry, together with two standard principles about grounding, yields the thesis.

Section~\ref{s:history} reviews the history of the question. Section~\ref{s:framework} sets out the framework, and Section~\ref{s:argument} the argument. Section~\ref{s:entanglement} adds a second, independent argument from quantum entanglement, which also covers the case in which some parts are simple. Section~\ref{s:levels} relates the result to MET-10, Section~\ref{s:cases} considers further cases, Section~\ref{s:literature} recent literature, and Section~\ref{s:objections} objections.

\section{A history of answers}\label{s:history}

\subsection{Parmenides and the Eleatics}

The earliest monism is existence monism. In the Way of Truth, Parmenides argues that what is is ungenerated and imperishable, whole, of one kind, unmoved and complete (DK 28 B8); there is no coming to be, because coming to be would require that what is come from what is not, and what is not cannot be spoken or thought. Melissus and Zeno defended the conclusion that what is is one and changeless. The argument, as MET-21 shows, equivocates between the existential and the predicative uses of \enquote{is}: from the premise that what is not does not exist, it infers that nothing can be other than it is. Once the equivocation is removed, the argument establishes neither that there is only one thing nor that there is no change. The case that defeats Eleatic monism is the ordinary case of change: a leaf that is green and then is not, without anything coming from what does not exist. Eleatic monism is not the view we defend, and nothing in our argument uses its premises.

\subsection{Spinoza}

Spinoza's monism is closer to priority monism. In the \emph{Ethics}, there is only one substance, God or nature, and except God no substance can be or be conceived (\emph{Ethics} I, P14). Finite things are modes of the one substance: they exist, but they exist in another and are conceived through another, namely the substance of which they are modifications. Spinoza's argument proceeds from the definition of substance as what exists in itself and is conceived through itself, and from the proposition that two substances cannot share an attribute (I, P5). The case that defeats the argument is the one Leibniz pressed: two substances might share an attribute and differ in their modes, so that P5 does not follow from the definitions. Spinoza's conclusion survives the failure of his argument, however, and the structure of his view, one fundamental whole and many derivative but real finite things, is the structure of the view we defend. Our whole is the concrete cosmos, not a being with infinitely many attributes, and our argument proceeds from the physics and metaphysics of composition rather than from the definition of substance.

\subsection{British idealism and the reality of relations}

The monism of the British idealists was argued from the nature of relations. \citet{Bradley1893} held that relations are unreal, because a relation that relates its terms must itself be related to them by further relations, without end; and he concluded that the many related things are appearances of a single, non-relational Absolute. The doctrine that all relations are internal, grounded in the natures of their terms, was meant to have the same consequence: if every relation of a thing to another is part of its nature, the nature of each thing includes its relations to everything, and nothing is what it is apart from the whole. \citet{Russell1903} and \citet{Moore1919} replied that some relations are external, holding independently of the natures of their terms, and that the argument from internal relations fails. MET-22 argues that Bradley's regress is not vicious and that relations are real, so that the idealist argument fails at its first step. The case that defeats it is any external relation, such as the distance between two bodies: it relates its terms without a further relating relation, and it is not part of the nature of either. Our argument does not depend on the doctrine of internal relations and is compatible with the reality of external ones.

\subsection{Pluralist orthodoxy}

The analytic tradition after Russell and Moore was pluralist in both senses. Russell's logical atomism held that the world consists of atomic facts about simples \citep{Russell1918}. \citet{Lewis1986Papers} held that all there is to the world is a vast mosaic of local matters of particular fact, the distribution of local qualities over the points of spacetime, and that everything else supervenes on it; Humean supervenience was for him a thesis about the priority of the local over the global. The case that defeats this form of pluralism is quantum entanglement. The state of an entangled pair is not fixed by the intrinsic states of its members together with their spatiotemporal relations: two pairs can agree on the latter and differ in the former \citep{Teller1986,Healey1991,Maudlin2007}. The whole has features that the local mosaic does not determine, and Lewis's thesis is, at least in the form he gave it, false of our world. We return to entanglement in Section~\ref{s:entanglement}.

\subsection{Existence monism revived}

\citet{HorganPotrc2008} defended a revived existence monism, which they called blobjectivism: there is exactly one concrete particular, the whole cosmos, and ordinary sentences about tables and persons are true by indirect correspondence to the way the blobject is, without there being tables or persons. Their main argument is that vagueness in ordinary objects cannot be accommodated in a correct ontology, so that the only concrete object without vague boundaries is the whole. The case that defeats the view is the existence of ordinary objects. MET-06 argues that tables and organisms exist, and that their parts compose them because they share a common fate; MET-18 argues that vague predicates have sharp but unknowable boundaries, so that the argument from vagueness does not require eliminating vague objects. Blobjectivism keeps the insight that the cosmos is special and pays for it with the denial of everything else; the present view keeps the insight without the price.

\subsection{Priority monism}

\citet{Schaffer2010} separated priority monism from existence monism and gave it its contemporary form. On his view, the cosmos is the one basic object, and its proper parts are derivative; he defended the view with three arguments. The argument from common sense holds that a whole is prior to its parts when it is integrated rather than a mere heap, and that the cosmos is integrated. The argument from quantum entanglement holds that the cosmos is in an entangled state, so that its properties are not fixed by those of its parts. The argument from gunk holds that there must be a fundamental basis, that the cosmos might be gunky, and that in a gunky world only the whole can be basic. Schaffer also imposed a tiling constraint, that the basic objects together compose the cosmos without overlap, which ensures that the choice between monism and pluralism is exhaustive. \citet{Schaffer2007} had earlier argued that the considerations which lead some to nihilism about composition lead, properly understood, to monism, and \citet{Schaffer2013} argued that the laws of nature describe the evolution of the cosmos as a whole.

Our argument differs from Schaffer's in three ways. We do not rely on the argument from common sense, since the integration of the cosmos is what is at issue. We do not rely on the argument from gunk in Schaffer's form, since it requires that there be a fundamental basis, which MET-10 denies is a requirement on grounding. And our main argument does not rely on entanglement, which we use only as a supplement. It relies instead on a feature of the composition of proper parts, their dependence on an environment and on the laws, that the cosmos does not share.

\section{Framework}\label{s:framework}

We speak of objects and of facts, and of the relation of partial grounding between facts, written $f \prec g$. A \emph{fact about} an object $x$ is a fact in which $x$ figures as a subject: that $x$ exists, that $x$ has some property, that $x$ stands in some relation, that $x$ is all there is. For each object $x$ we write $E(x)$ for the fact that $x$ exists.

\begin{definition}[Cosmos]\label{d:cosmos}
The \emph{cosmos} is the fusion of all concrete objects. It exists and is a composite, by the criterion of MET-06.
\end{definition}

\begin{definition}[Dependence; fundamental object]\label{d:fund}
An object $y$ \emph{depends on} an object $x$ if $x \neq y$ and some fact about $x$ partly grounds $E(y)$. An object is \emph{fundamental} if it depends on no object.
\end{definition}

\begin{definition}[Abstraction]\label{d:abstraction}
An object $y$ is an \emph{abstraction from} $x$ if $y$ is a proper part of $x$ and $y$ depends on $x$.
\end{definition}

\begin{definition}[Environment; totality fact]\label{d:env}
The \emph{environment} of a proper part $p$ of the cosmos is the fusion of the parts of the cosmos that do not overlap $p$. A \emph{totality fact} is a fact to the effect that some things are all there is, or that there is nothing of some kind beyond them; negative general facts, such as the fact that there is no $F$ that is not $G$, are totality facts.
\end{definition}

Definition~\ref{d:fund} follows \citet{Schaffer2010} in defining fundamentality for objects through the grounding of facts. Definition~\ref{d:abstraction} gives \enquote{abstraction} a precise sense: an abstraction from $x$ is a part of $x$ that exists in virtue of how $x$ is. The term is chosen because, on the view defended, to single out an ordinary object is to attend to one aspect of the cosmos, the portion of it in which a common fate obtains, and to abstract from the rest; but the definition carries no suggestion that the abstraction is performed by a mind, and abstractions so defined are concrete objects. Every proper part of the cosmos has an environment, by Definition~\ref{d:env}; the cosmos has none.

\begin{assumption}[Grounding]\label{a:ground}
Partial grounding is a strict partial order on facts: it is irreflexive and transitive.
\end{assumption}

\begin{assumption}[Necessitation]\label{a:nec}
If a set of facts fully grounds a fact $g$, then necessarily, if the members of the set obtain, $g$ obtains.
\end{assumption}

\begin{assumption}[No simples]\label{a:gunk}
Every proper part of the cosmos has proper parts.
\end{assumption}

Assumption~\ref{a:ground} is the assumption of MET-12-a and MET-10-a. As both note, it does not require grounding to be well-founded, and we shall not need well-foundedness. Assumption~\ref{a:nec} is standard \citep{Rosen2010,Fine2012}: a full ground is sufficient for what it grounds, and sufficiency in the metaphysical sense is necessitation. Assumption~\ref{a:gunk} is a consequence of the argument of MET-10-a (Section~4.1), on which the matter of the actual world is the portion of continuous fields in regions of spacetime, every region has proper subregions, and there is no level of facts at points; MET-10 concludes that every level of parts has a finer level. We use it only to ensure that every proper part is a composite, and Section~\ref{s:entanglement} shows that the thesis survives without it.

\begin{principle}[Existence first]\label{p:exist}
For every object $x$ and every fact $\varphi$ about $x$ other than $E(x)$, $E(x)$ partly grounds $\varphi$.
\end{principle}

\begin{principle}[Composition grounds]\label{p:comp}
If $p$ is a composite proper part of the cosmos and the $xx$ are its parts at some level, $E(p)$ is fully grounded in the fact that the $xx$ share a common fate under their relevant dynamics.
\end{principle}

\begin{principle}[Dynamics]\label{p:dyn}
For every plurality of proper parts of the cosmos, the fact that they share a common fate under their relevant dynamics, or fail to, is partly grounded in facts about which generalizations are laws.
\end{principle}

\begin{principle}[Cosmic totality]\label{p:tot}
Every totality fact is a fact about the cosmos.
\end{principle}

\ref{p:exist} is the grounding form of serious actualism, the view that nothing has properties or stands in relations without existing. It says more than that: not only is it necessary that whatever is some way exists, but its existing is part of what its being that way consists in. That Socrates is wise obtains in part because Socrates exists, and not conversely. The principle is accepted in some form by most theorists of ground who countenance facts about objects; it is the principle behind the familiar claim that a singular fact has its subject as a constituent and obtains only if its constituents do.

\ref{p:comp} is the grounding form of the criterion of MET-06. That criterion is a biconditional: the $xx$ compose something if and only if they share a common fate. \ref{p:comp} adds the direction of explanation for proper parts: a table exists because its parts are bound together under the perturbations of its environment, and not the other way round. The principle is restricted to proper parts, and the restriction is not ad hoc. For a proper part, the common-fate condition is a substantive condition on how the parts respond to perturbations from outside them. For the cosmos it is not: there are no admissible perturbations of the plurality of all things, since there is nothing outside it to produce them, and the condition is satisfied vacuously (MET-06-a, Section~6). A vacuous truth about the parts is not a fact about how they behave that could explain the existence of a whole; it holds because the parts are all there is, which is a totality fact. We return to the point in Corollary~\ref{c:composite}.

\ref{p:dyn} follows from Definition~3.2 of MET-06-a, on which the relevant dynamics of a plurality are the laws governing the interactions among its members, together with the regime of external influences to which its environment characteristically exposes it. Whether a push on one leg of a table carries the other legs with it depends on the laws of the interactions among the molecules of the table; had the laws been different, the same molecules in the same arrangement might not have shared a common fate. The regime of the environment likewise depends on the laws, since what the environment characteristically does is what it does under them.

\ref{p:tot} is the least familiar of our premises, and we defend it in Section~\ref{s:totality} below. It says that the subject of a fact to the effect that some things are all there is is the cosmos: to say that there is nothing beyond the $xx$ is to say that the cosmos is the fusion of the $xx$, and the second formulation shows what the fact is about.

\section{The argument}\label{s:argument}

\subsection{Laws and totality}

The first lemma connects the laws to the whole. MET-12 argues that laws are the axioms of the best system summarizing the mosaic, and MET-12-a (Proposition~3.2) shows that lawhood is global: whether a generalization is a law depends on the whole mosaic, including events far from any of its instances. Adding a single event anywhere can change which generalizations are laws.

\begin{lemma}[Laws need totality]\label{l:lawtot}
Every full ground of a fact about which generalizations are laws includes a totality fact.
\end{lemma}

\begin{proof}
Let $\lambda$ be the fact that some generalization $G$ is a law, and let $\Gamma$ be a full ground of $\lambda$. Suppose that $\Gamma$ includes no totality fact, so that it consists of facts about which events occur and what they are like, none of which says that these are all the events. By the globality of lawhood, for any set of such facts there is a possible mosaic that contains every event they describe together with further events, so many and so arranged that $G$ is not a law of the enlarged mosaic: for instance, a mosaic that adds counterinstances to $G$ in sufficient number that no true system containing $G$ is best. In that mosaic the members of $\Gamma$ obtain and $\lambda$ does not. This contradicts Assumption~\ref{a:nec}. Hence $\Gamma$ includes a totality fact. The same argument applies to the fact that a generalization is not a law, by adding instances that make it strong enough to enter the best system.
\end{proof}

Lemma~\ref{l:lawtot} makes explicit what the globality of lawhood implies. It is the analogue, for laws, of a point familiar from the theory of truthmakers: no collection of positive facts about particular things makes it true that all ravens are black, since the same facts would obtain in a world with one more raven, which is white; what is needed in addition is the fact that these are all the ravens \citep{Armstrong1997}. The best-system account inherits the point, since lawhood requires the truth of the axioms and the truth of a generalization requires a totality fact.

\subsection{The cosmic character of totality}\label{s:totality}

\ref{p:tot} says that totality facts are facts about the cosmos. The pluralist may grant Lemma~\ref{l:lawtot} and reply that the totality fact it requires is a fact about the many: that \emph{these} events, or \emph{these} particles, are all there are. That fact is about a plurality of proper parts, and it is not a fact about the cosmos.

We have three reasons to reject the reply. First, a plural totality fact says, of some things, that nothing exists beyond them; and the domain beyond which nothing exists is the cosmos. \enquote{The $xx$ are all there is} is necessarily equivalent to \enquote{the cosmos is the fusion of the $xx$}, and the second formulation identifies its subject: the extent of the cosmos is what the fact constrains. Second, by MET-10 there is no fundamental level, and every plurality of proper parts that are all there is is a level, a decomposition of the cosmos into proper parts. There is accordingly no privileged plurality whose totality fact the pluralist could take as basic: the totality fact at one level and the totality fact at any finer level say the same thing, that the cosmos contains nothing beyond what it contains, and they say it of different pluralities only because the cosmos can be decomposed in different ways. A fact that is invariant across every decomposition of a whole is a fact about the whole. Third, a totality fact is not grounded in the facts about the things it concerns taken one by one, or in any collection of them, by the reasoning of Lemma~\ref{l:lawtot}: each such collection could have obtained in a larger world. Totality facts are therefore not facts about parts in the sense in which a part's mass is; they are facts about how far reality extends, and reality, in the concrete case, is the cosmos.

We conclude that \ref{p:tot} is true. Let $T$ be the fact that the cosmos is all there is, that is, that there is no concrete object that is not part of it. Every totality fact about concrete things is either $T$ or partly grounded in it, since what concrete things there are beyond any given things is settled by the extent of the cosmos.

\subsection{The priority of the whole}

\begin{lemma}[Top-down]\label{l:topdown}
For every proper part $p$ of the cosmos, $T$ partly grounds $E(p)$.
\end{lemma}

\begin{proof}
By Assumption~\ref{a:gunk}, $p$ is a composite; let the $xx$ be its parts at some level. By \ref{p:comp}, $E(p)$ is fully grounded in the fact $\kappa$ that the $xx$ share a common fate under their relevant dynamics, so $\kappa \prec E(p)$. By \ref{p:dyn}, $\kappa$ is partly grounded in some fact $\lambda$ about which generalizations are laws. By MET-12, lawhood is grounded in the mosaic, so $\lambda$ has a full ground; by Lemma~\ref{l:lawtot}, every full ground of $\lambda$ includes a totality fact, which by \ref{p:tot} and the conclusion of Section~\ref{s:totality} is $T$ or is partly grounded in $T$. So $T \prec \lambda$, either directly or by transitivity. By transitivity again, $T \prec \lambda \prec \kappa \prec E(p)$ gives $T \prec E(p)$.
\end{proof}

\begin{lemma}[The cosmos depends on nothing]\label{l:nothing}
No fact about any proper part of the cosmos partly grounds $E(c)$, where $c$ is the cosmos.
\end{lemma}

\begin{proof}
Suppose that some fact $\psi$ about a proper part $p$ partly grounds $E(c)$. By \ref{p:exist}, either $\psi$ is $E(p)$ or $E(p) \prec \psi$; in either case, by transitivity, $E(p) \prec E(c)$. $T$ is a fact about $c$ other than $E(c)$, so by \ref{p:exist}, $E(c) \prec T$. By Lemma~\ref{l:topdown}, $T \prec E(p)$. By transitivity, $E(c) \prec E(c)$, which Assumption~\ref{a:ground} forbids.
\end{proof}

\begin{theorem}[Priority monism]\label{t:main}
The cosmos is fundamental, and it is the only fundamental concrete object.
\end{theorem}

\begin{proof}
Every concrete object other than the cosmos is a proper part of it, by Definition~\ref{d:cosmos}. By Lemma~\ref{l:nothing}, the cosmos depends on no proper part of itself, and so on no concrete object. It does not depend on any abstract object either: no fact about an abstract object alone could settle whether there is a concrete cosmos or how far it extends, and we take it, with MET-26, that the concrete world is not grounded in a mathematical or informational structure without something structured. So the cosmos is fundamental. For every proper part $p$, Lemma~\ref{l:topdown} shows that $T$, a fact about the cosmos, partly grounds $E(p)$, so $p$ depends on the cosmos and is not fundamental.
\end{proof}

\begin{corollary}[Abstraction]\label{c:abstraction}
Every ordinary object is an abstraction from the cosmos. In particular, tables and organisms exist and are abstractions from the cosmos.
\end{corollary}

\begin{proof}
By MET-06-a (Corollary~4.3), tables and organisms exist. Every ordinary object is a proper part of the cosmos and, by the proof of Theorem~\ref{t:main}, depends on it; by Definition~\ref{d:abstraction}, it is an abstraction from it.
\end{proof}

\begin{corollary}[The cosmos is a composite]\label{c:composite}
The proper parts of the cosmos at any level compose it, and the fact that they share a common fate is partly grounded in $E(c)$ and does not partly ground it.
\end{corollary}

\begin{proof}
Let the $xx$ be the members of a level. They share a common fate vacuously, since they have no environment and so no admissible perturbations (MET-06-a, Section~6); by the criterion of MET-06 they compose something, and what they compose is the cosmos. The fact that they share a common fate holds because there is nothing outside them, and so it is partly grounded in $T$, which by \ref{p:exist} is partly grounded in $E(c)$. By Assumption~\ref{a:ground}, a fact partly grounded in $E(c)$ does not partly ground $E(c)$.
\end{proof}

Corollary~\ref{c:composite} answers the question how the cosmos can be both a composite and fundamental. For a proper part, composition is the ground of existence: the table exists because its parts are bound under the perturbations of an environment. For the cosmos, the direction is reversed: its parts compose it because it is all there is, and their sharing a common fate is a consequence of its existence and extent, not a ground of them. The criterion of MET-06 applies to both, and the cosmos is a composite in exactly the sense in which a table is; what differs is the direction of grounding, which, as MET-06-a notes, the special composition question leaves open.

\begin{proposition}[Environment and dependence]\label{p:env}
Every concrete object that has an environment depends on the cosmos, and the cosmos is the only concrete object without an environment.
\end{proposition}

\begin{proof}
An object with an environment is a proper part of the cosmos, by Definition~\ref{d:env}, and so depends on the cosmos by the proof of Theorem~\ref{t:main}. Conversely, every proper part $p$ has an environment, since the cosmos is a composite with $p$ as a proper part, and by the supplementation principle of mereology some part of the cosmos does not overlap $p$.
\end{proof}

Proposition~\ref{p:env} states the asymmetry on which the argument rests in its most general form. It also explains why priority monism is a thesis about the cosmos in particular, rather than an instance of a general principle that wholes are prior to their parts. Tables are not prior to their legs: a table exists because its parts share a common fate, and its parts share a common fate partly because of their own interactions. The cosmos is prior to its parts not because it is a whole but because it is the only whole without an outside.

\section{The argument from entanglement}\label{s:entanglement}

The argument of Section~\ref{s:argument} uses Assumption~\ref{a:gunk} to ensure that every proper part is a composite, so that its existence is grounded in a composition fact. Suppose instead that there are simples. Their existence is not grounded in composition, and Lemma~\ref{l:topdown} does not apply to them. A second argument, from quantum entanglement, closes the gap; it is close to the argument of \citet{Schaffer2010}, and it applies to the natures of objects as well as to their existence.

\begin{definition}[Reduced state]\label{d:reduced}
If a composite system is in a state $\Psi$, the \emph{reduced state} of one of its subsystems is the state determined by $\Psi$ for the observables of that subsystem alone.
\end{definition}

\begin{assumption}[Non-separability]\label{a:entangle}
The cosmos is in a state $\Psi$ such that (i) the physical state of every proper part is its reduced state relative to $\Psi$, and (ii) $\Psi$ is not determined by the reduced states of the members of any level together with their spatiotemporal relations.
\end{assumption}

\begin{principle}[Determination]\label{p:det}
If the fact that $x$ is in state $\Psi$ necessitates the fact that $y$ is in state $\rho$, by an operation that yields $\rho$ from $\Psi$ and not conversely, and $y$ is a part of $x$, then the fact that $x$ is in $\Psi$ partly grounds the fact that $y$ is in $\rho$.
\end{principle}

Assumption~\ref{a:entangle} is standard quantum mechanics. Clause (i) says that the state of a subsystem is obtained from the state of the system by tracing out the rest; clause (ii) says that the global state is not obtained from the reduced states, which is what entanglement consists in \citep{Teller1986,Healey1991}. \citet{IsmaelSchaffer2020} argue that non-separability in this sense is common ground among the main interpretations of quantum mechanics. On the Everett interpretation, which SCI-09 defends, the state of the cosmos is the universal state, and every other physical state is a reduced state of it; the assumption then holds in the strongest form. \ref{p:det} says that when one state is obtained from another by a many-one operation, and not conversely, the one obtained is the derivative one; it is the principle by which we ordinarily take the reduced state of a subsystem to be a feature of the subsystem's embedding in a larger system.

\begin{proposition}[Natures are derivative]\label{p:natures}
Given Assumption~\ref{a:entangle} and \ref{p:det}, the fact that the cosmos is in state $\Psi$ partly grounds the fact that each proper part is in the state it is in, whether or not the part is simple.
\end{proposition}

\begin{proof}
By clause (i), the state of each proper part is its reduced state, which is obtained from $\Psi$ by tracing out the rest of the cosmos; the fact that the cosmos is in $\Psi$ therefore necessitates it. By clause (ii), $\Psi$ is not obtained from the reduced states of any level. \ref{p:det} applies and gives the result.
\end{proof}

Proposition~\ref{p:natures} concerns the natures of proper parts, not their existence. To extend it to existence, we need only the observation that a simple, if there are any, is individuated by its state: to be this electron rather than another is to be the subsystem whose reduced state is this one. Then the existence of a simple is grounded, at least in part, in the fact that the cosmos is in a state that has it as a subsystem. We do not rest the main argument on this observation, because it depends on contested views about the individuation of quantum systems, and because Assumption~\ref{a:gunk} makes it unnecessary. The point of the present section is that the monist does not need to bet on gunk: if there are no simples, Lemma~\ref{l:topdown} covers every proper part; if there are, Proposition~\ref{p:natures} shows that their natures, and plausibly their existence, are grounded in the state of the cosmos. A route to monism that does not appeal to entanglement at all is the subject of MET-09-b.

\section{Monism without a fundamental level}\label{s:levels}

MET-10 argues that there is no fundamental level of reality: every level is grounded in a lower one, and the descent is not vicious. Taken together with the present thesis, this may seem inconsistent. If nothing is fundamental, the cosmos is not; if the cosmos is fundamental, something is.

The inconsistency is only apparent, and resolving it shows how the two results fit. A level, in the sense of MET-10-a (Definition~3.1), is a plurality of non-overlapping proper parts of the cosmos that together compose it. The cosmos is not a proper part of itself, and so it is not a level and is not a member of any level. MET-10 says that no level is fundamental, that is, that every level of proper parts is grounded in a finer level of proper parts, without end. The present thesis says that the cosmos is fundamental and that every proper part depends on it. The descent concerns proper parts: there is no smallest level. The fundamentality concerns the whole: there is a greatest object, and it is the only one that depends on nothing.

The two groundings run in different directions, and both are real. A table's existence and character are grounded in facts about its parts and their relations, as MET-10 says, and its existence is also grounded, through the composition facts and the laws, in the totality fact $T$, as Lemma~\ref{l:topdown} says. Nothing in MET-10 requires that facts at a level be grounded \emph{only} in finer levels, and nothing here requires that facts about proper parts be grounded \emph{only} in the whole. The structure is the one exhibited in MET-10-a (Proposition~6.1): an infinite descending chain of facts at successively finer levels, each partly grounded in the next, and below them all, as a lower bound in the grounding order that is not a member of the chain, a single fact about the cosmos that partly grounds every one of them. Mereologically, the levels descend from the cosmos toward ever smaller parts; in the order of grounding, they descend toward the cosmos.

\begin{proposition}[Joint consistency]\label{p:joint}
The principles of the present manuscript are consistent with the conclusion of MET-10 and with Assumption~\ref{a:ground}.
\end{proposition}

\begin{proof}
Take the model of MET-10-a (Propositions~4.5 and~6.1): a cube of homogeneous matter, its decompositions into congruent sub-cubes of side $2^{-n}$, and facts at each level fully grounded by facts at the next. In place of the single fact about the whole in that model, take two facts, $E(c)$ and $T$. Let the grounding relation be the transitive closure of the following pairs: $E(c)$ partly grounds $T$ and every other fact about the cube as a whole; $T$ partly grounds every composition fact, including the vacuous common fate of the cubes at the coarsest level; the existence of a cube partly grounds every other fact about it; and every fact at level $n+1$ about the sub-cubes of a cube partly grounds the facts at level $n$ about that cube, among them its existence. Every generating pair either leads from $E(c)$ or $T$ into the levels or from a finer level to a coarser one, and none leads back, so the generating relation has no cycles and its transitive closure is a strict partial order. Every level is fully grounded by the next, as MET-10 requires; no fact about a cube grounds $E(c)$, so the cosmos depends on nothing; and $T$ partly grounds the existence of every cube, so every cube depends on the cosmos.
\end{proof}

The combination of priority monism with infinite descent was anticipated by Schaffer, for whom the possibility of gunk was a reason for monism. On his argument, the monist secures a fundamental basis in a gunky world where the pluralist cannot. We do not use the argument in that form, because, with MET-10, we deny that grounding must be well-founded in the sense of having no infinite descending chains; the descent among proper parts would be harmless even if nothing were fundamental. The monist's claim is not that a foundation is required and only the whole can supply it, but that the whole is in fact prior to its parts, for the reasons of Section~\ref{s:argument}. The combined view has the virtue Schaffer sought, that every fact is partly grounded in something fundamental, without making that virtue a requirement.

It may be asked why grounding should run upward among proper parts, from smaller to larger, and downward from the cosmos to everything else. The answer is Proposition~\ref{p:env}. Among proper parts, the existence of a whole is grounded in the binding of its parts under an environment, and the binding is a matter of the parts' interactions; so wholes are grounded in their parts. The cosmos is not bound under an environment, and the binding of its parts is a consequence of its being all there is; so it is not grounded in them. The two directions are not two relations: by MET-14, grounding is a single relation, and its apparent varieties are that relation applied at different levels. What differs between the cosmos and the table is not the relation but the facts it relates.

\section{Further cases}\label{s:cases}

\paragraph{Island universes.} Suppose the cosmos consists of two regions, $A$ and $B$, between which there is no causal interaction. One might think that each region is then a closed system, with no environment that perturbs it, so that the cosmos is not the only object without an outside, and the argument yields two fundamental objects. It does not. Each region has an environment in the sense of Definition~\ref{d:env}, namely the other region, although the environment exerts no perturbations. If the parts of $A$ share a common fate vacuously, they do so because $B$ does not interact with them, which is a fact about the laws and about what $B$ contains; and the laws, by Lemma~\ref{l:lawtot}, depend on the totality fact about the cosmos as a whole, including $B$. So $A$ depends on the cosmos, and so does $B$. The case shows that the argument does not depend on the cosmos being causally integrated, as the argument from common sense does; it depends only on the globality of the laws.

\paragraph{Junk.} \citet{Bohn2009} pointed out that a world might be junky: every object might be a proper part of a further object, so that there is no greatest object and no cosmos. In such a world, priority monism has no candidate for its one fundamental object. We claim only that the actual world contains a cosmos, which MET-06-a shows, since the plurality of all things shares a common fate vacuously and so composes something. We do not claim that monism is necessary, and the possibility of junk does not refute the actual truth of monism. The asymmetry between gunk and junk is real and principled: the criterion of MET-06 ensures that there is a greatest object, since the plurality of everything has no environment, while nothing ensures that there is a smallest one.

\paragraph{The four-dimensional cosmos.} We have stated the argument for the cosmos at a time, but the best-system account of laws concerns the whole history of the world, and Lemma~\ref{l:lawtot} requires the totality fact for the whole of spacetime. The fundamental object is therefore the four-dimensional cosmos, the fusion of everything at every time, and the objects that depend on it include not only ordinary objects at a time but their temporal parts and their careers. This fits the diachronic extension of the common-fate criterion in MET-06-a (Section~6), on which the temporal parts of a persisting object are those linked by the relevant dynamics.

\paragraph{Organisms.} Organisms are sometimes cited as wholes prior to their parts: the heart is what it is because of its role in the body, and a heart outside a body is a heart only by courtesy. Nothing in our argument requires that organisms be prior to their organs, and Proposition~\ref{p:env} explains why they need not be. An organism has an environment, and its existence is grounded in the binding of its parts under that environment, as MET-06 says. Whatever priority the organism has over its organs is a priority of function, or of the individuation of kinds of organ, not of existence. The case illustrates that the monist's whole is special because it has no outside, not because it is a whole.

\paragraph{The Everettian state.} On the Everett interpretation of quantum mechanics, which SCI-09 defends, the fundamental physical object is the universal state, whose branches and subsystems are structures within it. The interpretation is monist in spirit, and it gives Assumption~\ref{a:entangle} its strongest form. Our argument does not require it: Lemma~\ref{l:topdown} uses only the best-system account of laws and the criterion of composition, which are neutral among interpretations. But on the Everett interpretation, the result that ordinary objects are abstractions from the cosmos acquires a physical sense: a table is a stable, decohered structure in the universal state, and to single it out is to attend to that structure and abstract from the rest.

\section{Recent literature}\label{s:literature}

Since \citet{Schaffer2010}, priority monism has been examined from several directions. \citet{Sider2007} argued that the monist's ideology is costly: since the fundamental facts are all facts about one object, the monist must state them by attributing to the cosmos a vast stock of properties, which makes the statement of the laws less simple than it is for the pluralist. We answer the argument in Objection~\ref{o:ideology}. \citet{Cameron2010} argued for priority monism from truthmaker theory: on a Humean view, on which there are no necessary connections between distinct existences, the only entity that can make true the truths about what there is and what there is not is the world as a whole. His argument and ours share a reliance on totality: what makes it true that there is nothing more is not any part of the world. Our argument is not committed to truthmaker maximalism, however, and it reaches the conclusion through the laws and composition rather than through truthmaking in general.

\citet{Bohn2009} raised the possibility of junk, which we discussed in Section~\ref{s:cases}. \citet{IsmaelSchaffer2020} defended the argument from entanglement against the objection that it depends on a particular interpretation of quantum mechanics, arguing that non-separability is common ground. \citet{Schaffer2013} argued that the laws of nature describe the evolution of the cosmos as a whole, so that the cosmos is the one thing that acts. Our use of the laws is different: we do not claim that the laws govern the evolution of the whole, which MET-12 denies, but only that which generalizations are laws depends on the whole, which MET-12 asserts.

A more general challenge concerns grounding itself. \citet{Wilson2014} argued that there is no single relation of grounding that does the work theorists assign to it, but only a variety of specific dependence relations, such as composition, realization and determination. The challenge, if successful, would require us to restate the argument in terms of the specific relations involved: the dependence of the existence of composites on the dynamics of their parts, of lawhood on the mosaic, and of the reduced states on the global state. We think the restatement possible, but MET-14 argues that grounding is a single relation whose varieties are that relation at different levels, and we rely on its result. A different challenge, due to \citet{Barnes2018}, holds that dependence need not be asymmetric; we answer it in Objection~\ref{o:symmetric}.

Existence monism has continued to find defenders in the tradition of \citet{HorganPotrc2008}, who have argued that a minimal ontology with one concrete particular makes better sense of vagueness than any ontology of ordinary objects. Our disagreement with them is not about the special status of the cosmos, which we share, but about ordinary objects, which MET-06 and MET-18 show to be both real and sharply bounded.

\section{Objections and replies}\label{s:objections}

\begin{objection}[Symmetric dependence]\label{o:symmetric}
The argument for Lemma~\ref{l:nothing} uses the asymmetry of grounding to conclude that, since the parts depend on the cosmos, the cosmos does not depend on them. But dependence can be symmetric: there are cases in which each of two things depends on the other \citep{Barnes2018}. The cosmos may depend on its parts as they depend on it.
\end{objection}

\begin{reply}
The cases offered of symmetric dependence are cases of modal or identity dependence: each of two things could not exist without the other, or neither can be identified without reference to the other. Dependence in that sense can be symmetric, and nothing we say denies it; the cosmos and its parts may well be mutually dependent in the modal sense. Definition~\ref{d:fund} defines dependence by grounding, and grounding is a strict partial order (Assumption~\ref{a:ground}), as MET-12-a and MET-10-a also assume. Lemma~\ref{l:nothing} derives the asymmetry between the cosmos and its parts from that assumption together with \ref{p:exist}; it does not assume it. An objector who denies the asymmetry of grounding gives up the claim that grounding is explanatory, since a fact cannot be part of the explanation of what it is part of the explanation of.
\end{reply}

\begin{objection}[Ideology]\label{o:ideology}
If the cosmos is the only fundamental object, every fundamental fact is a fact about it, and the laws must be stated in terms of its properties: an enormous number of distributional properties, one for each way matter could be arranged. The pluralist states the laws in terms of a few properties of small things. Monism is therefore less simple \citep{Sider2007}.
\end{objection}

\begin{reply}
The objection assumes that the monist must state the fundamental facts in a vocabulary of properties of the cosmos alone. Our thesis is about the existence of objects, not about the vocabulary in which the facts are best stated. The best system, on MET-12, is stated in the vocabulary of genuine properties and relations, which MET-02 argues are few, and those properties are instantiated by proper parts at every level, as MET-10-a (Proposition~4.5) shows they can be. The monist can accept the pluralist's statement of the laws word for word. What she denies is that the parts mentioned in that statement exist independently of the whole: the laws are stated in terms of mass and charge at regions, and which regions host objects, and that these are all the regions, are facts about the cosmos. The ideological cost the objection identifies is the cost of a monist who insists on stating everything as a property of the one; it is not a cost of the thesis.
\end{reply}

\begin{objection}[Derivative means unreal]\label{o:unreal}
If ordinary objects are abstractions, they are not really there; to call a table an abstraction is to say that it is an aspect under which we consider the cosmos, and so to adopt existence monism in a more polite form.
\end{objection}

\begin{reply}
An abstraction, by Definition~\ref{d:abstraction}, is a proper part of the cosmos whose existence is grounded in a fact about the cosmos. Grounded facts are facts; derivative objects are objects. A table exists, by MET-06-a (Corollary~4.3) and Corollary~\ref{c:abstraction}, and it is a concrete part of the cosmos with mass, location and causal powers. That it exists in virtue of something does not make it exist less, any more than a conjunction is less true for being true in virtue of its conjuncts. The difference from existence monism is the difference between a derivative existent and a non-existent, and there is no politer form of the second.
\end{reply}

\begin{objection}[Plural totality]\label{o:plural}
The totality fact needed for the laws is the fact that \emph{these} things are all there is. That is a fact about many things, and the pluralist can take it as fundamental, alongside the facts about the things themselves. Then the laws, and the composition of ordinary objects, are grounded in facts about the many, and the cosmos is grounded in them too.
\end{objection}

\begin{reply}
We gave three reasons in Section~\ref{s:totality} for counting totality facts as facts about the cosmos. The reply here adds a dilemma. The pluralist's plural totality fact concerns either a privileged plurality, the members of a fundamental level, or every plurality that is all there is. The first horn is excluded by MET-10: there is no fundamental level, and so no privileged plurality. On the second horn, the pluralist has infinitely many totality facts, one for each level, which say the same thing of different pluralities and which stand in grounding relations to one another only through the composition facts relating the levels. What they share, and what makes them all true, is that the cosmos extends no further than it does. The pluralist who takes all of them as fundamental posits infinitely many fundamental facts where one suffices; the pluralist who takes one of them as fundamental must say which, and there is no answer.
\end{reply}

\begin{objection}[The Humean mosaic]\label{o:mosaic}
MET-12 is a Humean account of laws, and Humeanism is pluralist: the mosaic is a distribution of local qualities, and the laws supervene on it. The monist cannot use a pluralist theory of laws as a premise.
\end{objection}

\begin{reply}
The best-system account requires that laws be grounded in the mosaic, and MET-12-a states this as its Principle P1, on which the instances of a law are among the facts that ground its being a law. It does not require that the mosaic be grounded in its local constituents rather than conversely, and MET-12-a itself gives up Humean supervenience in its no-necessary-connections form. What our argument needs from MET-12 is Proposition~3.2 of MET-12-a, that lawhood is global, and Lemma~\ref{l:lawtot} shows that globality brings in a totality fact. A Humean who accepts the best-system account is therefore already committed to the dependence of lawhood on a fact about the whole, and our argument draws the consequence for composition.
\end{reply}

\section{What the result does not establish}\label{s:limits}

The result does not establish existence monism: tables and organisms exist (Corollary~\ref{c:abstraction}). It does not establish that monism is necessary; it is a thesis about the actual world, which contains a cosmos and whose laws are global, and a junky world or a world whose laws were local would not satisfy its premises. It does not establish that every whole is prior to its parts; on the contrary, Proposition~\ref{p:env} shows that ordinary wholes are not. It does not establish that there is a fundamental level, and it is consistent with the conclusion of MET-10 that there is none (Proposition~\ref{p:joint}). It does not depend on the argument from entanglement, which Section~\ref{s:entanglement} offers as a supplement, nor on any particular interpretation of quantum mechanics. Finally, it does not say what the intrinsic character of the cosmos is; on the view of MND-01, physics describes the structure of the world and not its intrinsic nature, and nothing here decides that question.

\section{Conclusion}

The cosmos is the only fundamental object. Every proper part of it has an environment, and whether its parts compose it depends on how they respond to that environment under the laws, which depend on the whole mosaic, including the fact that the cosmos contains nothing more. The cosmos has no environment, and its parts compose it because it is all there is. Given that grounding is a strict partial order and that facts about an object are grounded in part in its existence, the cosmos depends on nothing, and every ordinary object is an abstraction from it: a real, concrete, derivative part. The view is consistent with the composition of ordinary objects by common fate, with the existence of the cosmos as a composite, and with the absence of a fundamental level of parts.

\begin{thebibliography}{99}\small
\bibitem[Armstrong(1997)]{Armstrong1997} Armstrong, D.~M. (1997). \emph{A World of States of Affairs}. Cambridge University Press.
\bibitem[Barnes(2018)]{Barnes2018} Barnes, E. (2018). Symmetric dependence. In R.~Bliss and G.~Priest (eds.), \emph{Reality and Its Structure: Essays in Fundamentality}. Oxford University Press.
\bibitem[Bohn(2009)]{Bohn2009} Bohn, E.~D. (2009). Must there be a top level? \emph{The Philosophical Quarterly}, 59(235), 193--201.
\bibitem[Bradley(1893)]{Bradley1893} Bradley, F.~H. (1893). \emph{Appearance and Reality}. Swan Sonnenschein.
\bibitem[Cameron(2010)]{Cameron2010} Cameron, R.~P. (2010). From Humean truthmaker theory to priority monism. \emph{No\^us}, 44(1), 178--198.
\bibitem[Fine(2012)]{Fine2012} Fine, K. (2012). Guide to ground. In F.~Correia and B.~Schnieder (eds.), \emph{Metaphysical Grounding: Understanding the Structure of Reality}. Cambridge University Press.
\bibitem[Healey(1991)]{Healey1991} Healey, R.~A. (1991). Holism and nonseparability. \emph{The Journal of Philosophy}, 88(8), 393--421.
\bibitem[Horgan and Potr\v{c}(2008)]{HorganPotrc2008} Horgan, T. and Potr\v{c}, M. (2008). \emph{Austere Realism: Contextual Semantics Meets Minimal Ontology}. MIT Press.
\bibitem[Ismael and Schaffer(2020)]{IsmaelSchaffer2020} Ismael, J. and Schaffer, J. (2020). Quantum holism: Nonseparability as common ground. \emph{Synthese}, 197(10), 4131--4160.
\bibitem[Lewis(1986)]{Lewis1986Papers} Lewis, D. (1986). \emph{Philosophical Papers, Volume II}. Oxford University Press.
\bibitem[Maudlin(2007)]{Maudlin2007} Maudlin, T. (2007). \emph{The Metaphysics Within Physics}. Oxford University Press.
\bibitem[Moore(1919)]{Moore1919} Moore, G.~E. (1919). External and internal relations. \emph{Proceedings of the Aristotelian Society}, 20, 40--62.
\bibitem[Rosen(2010)]{Rosen2010} Rosen, G. (2010). Metaphysical dependence: Grounding and reduction. In B.~Hale and A.~Hoffmann (eds.), \emph{Modality: Metaphysics, Logic, and Epistemology}. Oxford University Press.
\bibitem[Russell(1903)]{Russell1903} Russell, B. (1903). \emph{The Principles of Mathematics}. Cambridge University Press.
\bibitem[Russell(1918)]{Russell1918} Russell, B. (1918--19). The philosophy of logical atomism. \emph{The Monist}, 28--29.
\bibitem[Schaffer(2007)]{Schaffer2007} Schaffer, J. (2007). From nihilism to monism. \emph{Australasian Journal of Philosophy}, 85(2), 175--191.
\bibitem[Schaffer(2010)]{Schaffer2010} Schaffer, J. (2010). Monism: The priority of the whole. \emph{The Philosophical Review}, 119(1), 31--76.
\bibitem[Schaffer(2013)]{Schaffer2013} Schaffer, J. (2013). The action of the whole. \emph{Aristotelian Society Supplementary Volume}, 87(1), 67--87.
\bibitem[Sider(2007)]{Sider2007} Sider, T. (2007). Against monism. \emph{Analysis}, 67(1), 1--7.
\bibitem[Teller(1986)]{Teller1986} Teller, P. (1986). Relational holism and quantum mechanics. \emph{The British Journal for the Philosophy of Science}, 37(1), 71--81.
\bibitem[Wilson(2014)]{Wilson2014} Wilson, J.~M. (2014). No work for a theory of grounding. \emph{Inquiry}, 57(5--6), 535--579.
\end{thebibliography}

\end{document}
